\documentclass[11pt,a4paper]{article}
\usepackage[margin=1in]{geometry}
\usepackage[T1]{fontenc}
\usepackage[utf8]{inputenc}
\usepackage{lmodern}
\usepackage{hyperref}
\usepackage{enumitem}
\usepackage{xcolor}
\hypersetup{
  colorlinks=true,
  linkcolor=blue!60!black,
  urlcolor=blue!60!black,
  pdftitle={Crack the Code End User Guide},
  pdfauthor={Crack the Code Team}
}

\setlist[itemize]{leftmargin=1.6em}
\setlist[enumerate]{leftmargin=1.8em}

\title{\textbf{Crack the Code}\\End User Guide\\v1.2.1}
\author{National 5 Computing Science Worked Examples}
\date{\today}

\begin{document}
\maketitle
\tableofcontents
\newpage

\section{Purpose of this document}
This guide explains how students and teachers should use the Crack the Code web app.
It is written for end users and client stakeholders, and can be used as a reference during normal classroom use.

\section{System overview}
Crack the Code is an interactive learning platform for National 5 Computing Science.
It includes:
\begin{itemize}
  \item scaffolded worked examples,
  \item prediction, tracing, modify, and make tasks,
  \item adaptive hints and rich feedback,
  \item progress tracking and completion badges,
  \item teacher analytics and class management tools.
\end{itemize}

\section{Accessing the system}
\subsection{System links}
\begin{itemize}
  \item GitHub repository: \url{https://github.com/aaronhughes05/n5-programming-worked-examples}
  \item Live hosted site (Render): \url{https://n5-programming-worked-examples.onrender.com/}
  \item Admin view (Render): \url{https://n5-programming-worked-examples.onrender.com/admin}
\end{itemize}

\subsection{Web address}
Open the deployed site in a web browser:
\begin{itemize}
  \item \url{https://n5-programming-worked-examples.onrender.com}
\end{itemize}

\subsection{Sign in}
\begin{enumerate}
  \item Open the website.
  \item Enter your username and password on the sign-in screen.
  \item Select \textbf{Login}.
\end{enumerate}

\textbf{Note:} Users must be signed in to access learning pages.

\section{Navigation}
The app bar provides direct links to:
\begin{itemize}
  \item Home
  \item Examples
  \item Worksheets
  \item Final Assessment
\end{itemize}

The avatar menu (top right) provides account actions:
\begin{itemize}
  \item Change password
  \item Teacher mode (teacher accounts only)
  \item Sign out
\end{itemize}

\section{Student guide}
\subsection{Complete an example}
\begin{enumerate}
  \item Open an example from \textbf{Examples}.
  \item Read the \textbf{Problem Description} and \textbf{Learning Objectives}.
  \item Complete each step in order.
  \item Use \textbf{Check} buttons to validate answers.
  \item Continue until the completion screen appears.
\end{enumerate}

\subsection{How checkpoint gating works}
Some steps require earlier checkpoints to be correct before progression.
If a step does not unlock:
\begin{itemize}
  \item re-check missing answers,
  \item review feedback and hints,
  \item retry the task.
\end{itemize}

\subsection{Using hints and feedback}
\begin{itemize}
  \item Feedback appears after each check attempt.
  \item Hints can be revealed gradually to support learning.
  \item Worked hints are available for harder checkpoints.
\end{itemize}

\subsection{Final assessment}
The assessment combines all core concepts:
\begin{itemize}
  \item validation,
  \item loops and totals,
  \item list storage and traversal.
\end{itemize}

Students should complete examples first to improve readiness.

\subsection{Resume and persistence}
Progress is saved automatically to your account.
You can return later and continue from your saved state.

\section{Teacher guide}
\subsection{Open teacher mode}
\begin{enumerate}
  \item Sign in with a teacher account.
  \item Open the avatar menu.
  \item Select \textbf{Teacher Mode}.
\end{enumerate}

\subsection{Class and roster management}
In teacher mode you can:
\begin{itemize}
  \item create classes,
  \item add students to classes,
  \item remove students from classes,
  \item delete classes.
\end{itemize}

\subsection{Analytics available}
Teacher mode provides:
\begin{itemize}
  \item class progress summary (complete, in progress, not started),
  \item hint usage totals,
  \item checkpoint attempt analytics,
  \item most-missed checkpoints,
  \item per-student analytics via clickable student chips.
\end{itemize}

\subsection{Data export}
Teachers can export analytics as:
\begin{itemize}
  \item CSV
  \item JSON
\end{itemize}

\subsection{Demo data}
Teacher mode includes a seed action for inserting test/demo data.
Use this for presentations and controlled testing.

\section{Changing your password}
All authenticated users can change their own password.

\begin{enumerate}
  \item Open the avatar menu.
  \item Select \textbf{Change password}.
  \item Enter current password, new password, and confirmation.
  \item Select \textbf{Update password}.
\end{enumerate}

Password rules:
\begin{itemize}
  \item new password must differ from current password,
  \item new and confirmation must match,
  \item minimum strength/length is enforced by server policy.
\end{itemize}

\section{Troubleshooting}
\subsection{Cannot sign in}
\begin{itemize}
  \item Check username/password spelling.
  \item Confirm account exists.
  \item Contact teacher/admin if credentials are unknown.
\end{itemize}

\subsection{Step does not unlock}
\begin{itemize}
  \item Re-check required checkpoints in the previous step.
  \item Read rich feedback carefully and retry.
  \item Use hints where available.
\end{itemize}

\subsection{Teacher page access denied}
\begin{itemize}
  \item Only users with teacher role can access teacher mode.
  \item Sign in with a teacher account.
\end{itemize}

\subsection{Password change fails}
\begin{itemize}
  \item Ensure current password is correct.
  \item Ensure new password meets policy rules.
  \item If too many failed attempts occur, wait and retry.
\end{itemize}

\section{Good practice}
\begin{itemize}
  \item Students: complete each example before assessment.
  \item Teachers: review most-missed checkpoints for targeted support.
  \item All users: sign out after use on shared devices.
\end{itemize}

\section{Document control}
\begin{itemize}
  \item Artefact: Crack the Code end user guide
  \item Version: v1.2.1
  \item Intended audience: students, teachers, client stakeholders
\end{itemize}

\end{document}
